\documentclass[reqno]{amsart}
\usepackage[T1]{fontenc}
\usepackage{amsmath,amssymb,amsthm}
\usepackage{xr-hyper}
\usepackage[hidelinks]{hyperref}
\externaldocument[dt-]{20261125T002000--definite-theta-lifts__workbook}
\externaldocument[sw-]{20261125T015500--siegel-weil-in-the-smallest-case__workbook}
\externaldocument[dn-]{20261125T002100--definite-theta-lifts-and-norms__workbook}

% Chapter number; \ref{chap} gives it.
\newcounter{chap}
% Chapter numbers, generated from 20261125T012000--grh-workbook__guide.org
% by scripts/make-chapters.py.  Do not edit: change the order of the
% index and rerun the script (./preview-site.sh and CI do this).
%
% Usage in a sheet, after \newcounter{chap}:
%   % Chapter numbers, generated from 20261125T012000--grh-workbook__guide.org
% by scripts/make-chapters.py.  Do not edit: change the order of the
% index and rerun the script (./preview-site.sh and CI do this).
%
% Usage in a sheet, after \newcounter{chap}:
%   % Chapter numbers, generated from 20261125T012000--grh-workbook__guide.org
% by scripts/make-chapters.py.  Do not edit: change the order of the
% index and rerun the script (./preview-site.sh and CI do this).
%
% Usage in a sheet, after \newcounter{chap}:
%   \input{chapters}\chapnumber{<sheet ID>}
\makeatletter
\providecommand\chapters@set[2]{\@namedef{chapters@#1}{#2}}
\chapters@set{20261125T001000}{1}% zeros-to-ramp-sums
\chapters@set{20261125T001100}{2}% inert-primes-and-estimate-b
\chapters@set{20261125T005000}{3}% galois-fourier-and-averaging
\chapters@set{20261125T005500}{4}% columns-from-a-sphere
\chapters@set{20261125T000000}{5}% slope-method
\chapters@set{20261126T030100}{6}% a-twist-for-each-prime
\chapters@set{20261125T003000}{7}% the-chain-and-the-square-root
\chapters@set{20261125T010000}{8}% rank-lemma-and-zero-estimate
\chapters@set{20261125T010500}{9}% poincare-series-and-near-orthogonality
\chapters@set{20261124T223000}{10}% elliptic-curves-for-the-proof
\chapters@set{20261125T009500}{11}% fock-model-and-theta-intertwiners
\chapters@set{20261124T223100}{12}% complex-multiplication-and-inert-reduction
\chapters@set{20261125T008000}{13}% one-cm-determinant
\chapters@set{20261124T214500}{14}% supersingular-theta-divisibility
\chapters@set{20261124T214600}{15}% the-isotropic-plane
\chapters@set{20261125T013700}{16}% pairings-on-a-formal-group-of-height-two
\chapters@set{20261125T013500}{17}% hecke-nondegeneracy-on-the-tree
\chapters@set{20261126T030200}{18}% the-bound-is-attained
\chapters@set{20261125T009000}{19}% kronecker-limit-and-chowla-selberg
\chapters@set{20261125T002000}{20}% definite-theta-lifts
\chapters@set{20261125T015500}{21}% siegel-weil-in-the-smallest-case
\chapters@set{20261125T002100}{22}% definite-theta-lifts-and-norms
\chapters@set{20261125T016500}{23}% hecke-operators-and-theta-lifts
\chapters@set{20261125T016700}{24}% divisibility-along-the-chain
\chapters@set{20261125T004000}{25}% denominators-and-sha
\chapters@set{20261125T004100}{26}% tate-twists-and-sha
\chapters@set{20261126T020430}{27}% the-minus-part-by-gauss-sums
\chapters@set{20261125T014100}{28}% congruence-modules
\chapters@set{20261125T014500}{29}% eisenstein-congruences-and-selmer-classes
\chapters@set{20261125T014700}{30}% pseudocharacters-and-the-length-bound
\chapters@set{20261125T011000}{31}% galois-representations-black-box
\chapters@set{20261126T030300}{32}% the-data
\chapters@set{20261126T020210}{33}% bessel-poincare-series-and-their-gram-operators
\chapters@set{20261126T030400}{34}% the-seesaw-identity
\chapters@set{20261126T030500}{35}% divisibility-of-the-composite-lift
\chapters@set{20261126T030600}{36}% theorem-a-short-forms
\chapters@set{20261126T030700}{37}% the-determinant-of-a-chain
\chapters@set{20261126T030800}{38}% the-cost-of-one-lift
\chapters@set{20261126T030900}{39}% denominators-of-projections-to-lifts
\chapters@set{20261126T031000}{40}% twists-with-few-congruences
\chapters@set{20261126T031100}{41}% theorem-c-and-the-main-inequality
\providecommand\chapnumber[1]{%
  \@ifundefined{chapters@#1}%
    {\PackageError{chapters}{Sheet #1 is not listed in the index}%
      {List it in 20261125T012000--grh-workbook__guide.org\MessageBreak
       and run scripts/make-chapters.py.}}%
    {\edef\thechap{\@nameuse{chapters@#1}}}}
\makeatother
\chapnumber{<sheet ID>}
\makeatletter
\providecommand\chapters@set[2]{\@namedef{chapters@#1}{#2}}
\chapters@set{20261125T001000}{1}% zeros-to-ramp-sums
\chapters@set{20261125T001100}{2}% inert-primes-and-estimate-b
\chapters@set{20261125T005000}{3}% galois-fourier-and-averaging
\chapters@set{20261125T005500}{4}% columns-from-a-sphere
\chapters@set{20261125T000000}{5}% slope-method
\chapters@set{20261126T030100}{6}% a-twist-for-each-prime
\chapters@set{20261125T003000}{7}% the-chain-and-the-square-root
\chapters@set{20261125T010000}{8}% rank-lemma-and-zero-estimate
\chapters@set{20261125T010500}{9}% poincare-series-and-near-orthogonality
\chapters@set{20261124T223000}{10}% elliptic-curves-for-the-proof
\chapters@set{20261125T009500}{11}% fock-model-and-theta-intertwiners
\chapters@set{20261124T223100}{12}% complex-multiplication-and-inert-reduction
\chapters@set{20261125T008000}{13}% one-cm-determinant
\chapters@set{20261124T214500}{14}% supersingular-theta-divisibility
\chapters@set{20261124T214600}{15}% the-isotropic-plane
\chapters@set{20261125T013700}{16}% pairings-on-a-formal-group-of-height-two
\chapters@set{20261125T013500}{17}% hecke-nondegeneracy-on-the-tree
\chapters@set{20261126T030200}{18}% the-bound-is-attained
\chapters@set{20261125T009000}{19}% kronecker-limit-and-chowla-selberg
\chapters@set{20261125T002000}{20}% definite-theta-lifts
\chapters@set{20261125T015500}{21}% siegel-weil-in-the-smallest-case
\chapters@set{20261125T002100}{22}% definite-theta-lifts-and-norms
\chapters@set{20261125T016500}{23}% hecke-operators-and-theta-lifts
\chapters@set{20261125T016700}{24}% divisibility-along-the-chain
\chapters@set{20261125T004000}{25}% denominators-and-sha
\chapters@set{20261125T004100}{26}% tate-twists-and-sha
\chapters@set{20261126T020430}{27}% the-minus-part-by-gauss-sums
\chapters@set{20261125T014100}{28}% congruence-modules
\chapters@set{20261125T014500}{29}% eisenstein-congruences-and-selmer-classes
\chapters@set{20261125T014700}{30}% pseudocharacters-and-the-length-bound
\chapters@set{20261125T011000}{31}% galois-representations-black-box
\chapters@set{20261126T030300}{32}% the-data
\chapters@set{20261126T020210}{33}% bessel-poincare-series-and-their-gram-operators
\chapters@set{20261126T030400}{34}% the-seesaw-identity
\chapters@set{20261126T030500}{35}% divisibility-of-the-composite-lift
\chapters@set{20261126T030600}{36}% theorem-a-short-forms
\chapters@set{20261126T030700}{37}% the-determinant-of-a-chain
\chapters@set{20261126T030800}{38}% the-cost-of-one-lift
\chapters@set{20261126T030900}{39}% denominators-of-projections-to-lifts
\chapters@set{20261126T031000}{40}% twists-with-few-congruences
\chapters@set{20261126T031100}{41}% theorem-c-and-the-main-inequality
\providecommand\chapnumber[1]{%
  \@ifundefined{chapters@#1}%
    {\PackageError{chapters}{Sheet #1 is not listed in the index}%
      {List it in 20261125T012000--grh-workbook__guide.org\MessageBreak
       and run scripts/make-chapters.py.}}%
    {\edef\thechap{\@nameuse{chapters@#1}}}}
\makeatother
\chapnumber{<sheet ID>}
\makeatletter
\providecommand\chapters@set[2]{\@namedef{chapters@#1}{#2}}
\chapters@set{20261125T001000}{1}% zeros-to-ramp-sums
\chapters@set{20261125T001100}{2}% inert-primes-and-estimate-b
\chapters@set{20261125T005000}{3}% galois-fourier-and-averaging
\chapters@set{20261125T005500}{4}% columns-from-a-sphere
\chapters@set{20261125T000000}{5}% slope-method
\chapters@set{20261126T030100}{6}% a-twist-for-each-prime
\chapters@set{20261125T003000}{7}% the-chain-and-the-square-root
\chapters@set{20261125T010000}{8}% rank-lemma-and-zero-estimate
\chapters@set{20261125T010500}{9}% poincare-series-and-near-orthogonality
\chapters@set{20261124T223000}{10}% elliptic-curves-for-the-proof
\chapters@set{20261125T009500}{11}% fock-model-and-theta-intertwiners
\chapters@set{20261124T223100}{12}% complex-multiplication-and-inert-reduction
\chapters@set{20261125T008000}{13}% one-cm-determinant
\chapters@set{20261124T214500}{14}% supersingular-theta-divisibility
\chapters@set{20261124T214600}{15}% the-isotropic-plane
\chapters@set{20261125T013700}{16}% pairings-on-a-formal-group-of-height-two
\chapters@set{20261125T013500}{17}% hecke-nondegeneracy-on-the-tree
\chapters@set{20261126T030200}{18}% the-bound-is-attained
\chapters@set{20261125T009000}{19}% kronecker-limit-and-chowla-selberg
\chapters@set{20261125T002000}{20}% definite-theta-lifts
\chapters@set{20261125T015500}{21}% siegel-weil-in-the-smallest-case
\chapters@set{20261125T002100}{22}% definite-theta-lifts-and-norms
\chapters@set{20261125T016500}{23}% hecke-operators-and-theta-lifts
\chapters@set{20261125T016700}{24}% divisibility-along-the-chain
\chapters@set{20261125T004000}{25}% denominators-and-sha
\chapters@set{20261125T004100}{26}% tate-twists-and-sha
\chapters@set{20261126T020430}{27}% the-minus-part-by-gauss-sums
\chapters@set{20261125T014100}{28}% congruence-modules
\chapters@set{20261125T014500}{29}% eisenstein-congruences-and-selmer-classes
\chapters@set{20261125T014700}{30}% pseudocharacters-and-the-length-bound
\chapters@set{20261125T011000}{31}% galois-representations-black-box
\chapters@set{20261126T030300}{32}% the-data
\chapters@set{20261126T020210}{33}% bessel-poincare-series-and-their-gram-operators
\chapters@set{20261126T030400}{34}% the-seesaw-identity
\chapters@set{20261126T030500}{35}% divisibility-of-the-composite-lift
\chapters@set{20261126T030600}{36}% theorem-a-short-forms
\chapters@set{20261126T030700}{37}% the-determinant-of-a-chain
\chapters@set{20261126T030800}{38}% the-cost-of-one-lift
\chapters@set{20261126T030900}{39}% denominators-of-projections-to-lifts
\chapters@set{20261126T031000}{40}% twists-with-few-congruences
\chapters@set{20261126T031100}{41}% theorem-c-and-the-main-inequality
\providecommand\chapnumber[1]{%
  \@ifundefined{chapters@#1}%
    {\PackageError{chapters}{Sheet #1 is not listed in the index}%
      {List it in 20261125T012000--grh-workbook__guide.org\MessageBreak
       and run scripts/make-chapters.py.}}%
    {\edef\thechap{\@nameuse{chapters@#1}}}}
\makeatother
\chapnumber{20261125T016500}
\renewcommand{\thesection}{\thechap.\arabic{section}}

\theoremstyle{definition}
\newtheorem{problem}{Problem}[section]
\theoremstyle{plain}
\newtheorem{theorem}{Theorem}[section]

\newcommand{\Z}{\mathbb{Z}}
\newcommand{\Q}{\mathbb{Q}}
\newcommand{\F}{\mathbb{F}}
\newcommand{\C}{\mathbb{C}}
\newcommand{\ii}{\mathrm{i}}
\newcommand{\Oc}{\mathcal{O}}
\newcommand{\Cl}{\mathrm{Cl}}
\newcommand{\Nb}{\mathcal{N}}
\newcommand{\qbinom}[3]{\genfrac{[}{]}{0pt}{}{#1}{#2}_{#3}}

\setlength{\emergencystretch}{2em}

\title{Hecke operators and theta lifts}
\author{}
\date{November 25, 2026}

\begin{document}
\maketitle
\refstepcounter{chap}\label{chap}

\paragraph{Prerequisites.}
Even unimodular lattices, \(E_8\), theta series and the Hecke operators \(T_n\) on modular forms of level one \cite[Ch.~V and VII]{SeA}.  Quadratic forms over \(\F_p\) and the number of their isotropic vectors.  Gaussian binomial coefficients \(\qbinom{a}{b}{Q}\), which count the \(b\)-dimensional subspaces of \(\F_Q^a\) \cite[Ch.~3]{Andr}.  From Chapter~\ref{dt-chap}, \emph{Masses and the count of forms}: Section~\ref{dt-sec:lattices}, in particular Problems~\ref{dt-pr:real}, \ref{dt-pr:aut} and~\ref{dt-pr:rank-five}.  From Chapter~\ref{sw-chap}, \emph{The Siegel--Weil formula in the smallest case}: Problems~\ref{sw-pr:rank16} and~\ref{sw-pr:cusp5}.  From Chapter~\ref{dn-chap}, \emph{Definite theta lifts and their norms}: Section~\ref{dn-sec:kernel}.  For Section~\ref{sec:brandt}, maximal orders in a definite quaternion algebra over \(\Q\) \cite{Gr}, \cite{Pi}.

A theta lift is a sum over classes of lattices.  A Hecke operator on a class set is also a sum over lattices, namely over the neighbours of a lattice.  How do the two interact?  Eichler found that a theta series summed over the neighbours of a lattice is a Hecke transform of the theta series of the lattice, up to an explicit multiple of the identity \cite{Ei52}, \cite{Ei73}; see \cite{An} for Siegel modular forms and \cite{Ra} for the general theta correspondence.  We prove this in degree one, and check it on class sets with two elements, where the Hecke operator is a \(2 \times 2\) matrix.


\section{Neighbours}\label{sec:neighbours}

How does one move from a lattice to the other lattices of its genus?  Kneser's answer is to change it at one prime \cite{Kn}.

Let \(L\) be an even unimodular lattice of rank \(n = 2k\), with \(8 \mid n\), and let \(p\) be a prime.  A \emph{\(p\)-neighbour} of \(L\) is an even unimodular lattice \(M \subset L \otimes \Q\) with \([L : L \cap M] = [M : L \cap M] = p\).  Let \(\Nb_p(L)\) be the set of \(p\)-neighbours.

\begin{problem}\label{pr:nb-line}
Let \(p\) be odd.  For \(v \in L \setminus pL\) with \(v \cdot v \equiv 0 \pmod{2p^2}\), let \(L_v := \{x \in L : x \cdot v \equiv 0 \pmod p\}\).  Then \(M_v := L_v + \Z\, v/p\) is a \(p\)-neighbour of \(L\), and \(M_v\) depends only on the line \(\F_p \bar v \subset L/pL\).  Every \(p\)-neighbour arises in this way, from exactly one isotropic line of \(L/pL\).
\end{problem}

\begin{problem}\label{pr:nb-count}
The quadratic space \(L/pL\) over \(\F_p\) is split, so that
\[
  \lvert \Nb_p(L) \rvert = \frac{(p^k - 1)(p^{k-1} + 1)}{p - 1} .
\]
For \(L = E_8\) and \(p = 3, 5\) this is \(1120\) and \(19656\).
\end{problem}

For \(p = 2\) the same count holds, if \(v\) is chosen with \(v \cdot v \equiv 0 \pmod 8\) \cite{CL}; for \(E_8\) it is \(135\).  Since \(E_8\) is the only even unimodular lattice of rank \(8\), all its neighbours are copies of \(E_8\).  For \(n \ge 3\), Kneser proved that any two lattices of one genus are joined by a chain of \(p\)-neighbours \cite{Kn}.

Let \(\Cl\) be a class set, with \(\Gamma_L := \mathrm{O}(L)\).  The \emph{neighbour operator} on functions \(\varphi \colon \Cl \to \C\) is \((T_p\varphi)(L) := \sum_{M \in \Nb_p(L)} \varphi(M)\), and its matrix \(N_p(L, M) := \#\{M' \in \Nb_p(L) : M' \cong M\}\) is the \emph{Brandt matrix} of the class set at \(p\).  Let \(\langle \varphi, \psi \rangle := \sum_{L \in \Cl} \varphi(L)\overline{\psi(L)}/\lvert\Gamma_L\rvert\), as in Chapter~\ref{dt-chap}.

\begin{problem}\label{pr:nb-adjoint}
\(N_p(L, M)/\lvert\Gamma_L\rvert = N_p(M, L)/\lvert\Gamma_M\rvert\).  Hence \(T_p\) is self-adjoint for \(\langle\,,\rangle\), and the constant function is an eigenvector with eigenvalue \(\lvert\Nb_p(L)\rvert\).
\end{problem}

In rank \(16\) the class set is \(\{E_8 \perp E_8,\, D_{16}^+\}\), with \(\lvert\Gamma\rvert = 2 \cdot 696729600^2\) and \(2^{15}\cdot 16!\) \cite{CoS}.  By a machine enumeration of the \(32895\) two-neighbours of \(E_8 \perp E_8\), sorted by whether their roots span the lattice, \(14670\) of them are isometric to \(E_8 \perp E_8\) and \(18225\) to \(D_{16}^+\).

\begin{problem}\label{pr:rank16}
The Brandt matrix of the class set of rank \(16\) at \(p = 2\) is
\(\begin{pmatrix} 14670 & 18225 \\ 12870 & 20025 \end{pmatrix}\), with eigenvalues \(32895\) and \(1800\).
\end{problem}

The eigenvalue \(1800\) is \(15\,\tau(2) + (2^4 + 2^5 + \dots + 2^{10}) + 2^7\), where \(\tau(2) = -24\) is the coefficient of Ramanujan's \(\Delta\).  Chenevier and Lannes explain such identities by the Arthur parameters of the eigenvectors \cite{CL}.  In Section~\ref{sec:eichler} we show that the theta series of degree one cannot see this eigenvector.


\section{Eichler's commutation relation in degree one}\label{sec:eichler}

For \(L\) as in Section~\ref{sec:neighbours}, let \(r_L(m) := \#\{x \in L : x \cdot x = 2m\}\), with \(r_L(m) := 0\) unless \(m \in \Z_{\ge 0}\), and let \(\theta_L := \sum_m r_L(m)\, q^m\).  It is a modular form of weight \(k\) for \(\mathrm{SL}_2(\Z)\) \cite[Ch.~VII]{SeA}, and \(T_{p^2}\) acts on its coefficients by
\[
  a(m) \mapsto a(p^2m) + p^{k-1} a(m)\,[p \mid m] + p^{2k-2} a(m/p^2) .
\]
We compute \(\sum_{M \in \Nb_p(L)} r_M(m)\) by counting, for each vector, the neighbours that contain it.  Let \(p\) be odd.

\begin{problem}\label{pr:count-in}
Let \(x \in L \setminus pL\) with \(x \cdot x = 2m\).  The number of \(M \in \Nb_p(L)\) with \(x \in M\) is \((p^{2k-2} - 1)/(p - 1)\) if \(p \nmid m\), and \(1 + p(p^{k-1} - 1)(p^{k-2} + 1)/(p - 1)\) if \(p \mid m\).
\end{problem}

\begin{problem}\label{pr:count-out}
The vectors of \(\bigcup_{M \in \Nb_p(L)} M \setminus L\) are the \(y/p\) with \(y \in L \setminus pL\) and \(y \cdot y \equiv 0 \pmod{2p^2}\), and each lies in exactly one neighbour.
\end{problem}

Problems~\ref{pr:count-in} and~\ref{pr:count-out}, together with \(x \in pL \subset M\) for every \(M\), give the following.

\begin{theorem}[Eichler]\label{thm:eichler}
For odd \(p\),
\(\displaystyle \sum_{M \in \Nb_p(L)} \theta_M = T_{p^2}\theta_L + \frac{p^{2k-2} - p}{p - 1}\,\theta_L\).
\end{theorem}

\begin{problem}\label{pr:eichler-check}
Theorem~\ref{thm:eichler} for \(L = E_8\) and \(p = 3\) reads \(1120 = 757 + 363\).  The same identity of constant terms holds for \(p = 2\): \(135 = 73 + 62\).
\end{problem}

Let \(\theta(\varphi) := \sum_{L \in \Cl} \varphi(L)\, \theta_L / \lvert\Gamma_L\rvert\) be the \emph{theta lift} of a function on the class set.  It is the lift of Chapter~\ref{dn-chap} with the second lattice replaced by a vector of norm \(m\), the norms collected into a \(q\)-series.

\begin{problem}\label{pr:commute}
\(\theta(T_p\varphi) = \bigl(T_{p^2} + \frac{p^{2k-2} - p}{p-1}\bigr)\theta(\varphi)\).  If \(\theta(\varphi) \ne 0\) and \(T_p\varphi = \mu\varphi\), and \(\theta(\varphi)\) is an eigenform with \(T_p\theta(\varphi) = \lambda\theta(\varphi)\), then \(\mu = \lambda^2 - p^{k-1} + (p^{2k-2} - p)/(p-1)\).
\end{problem}

So the theta lift intertwines \(T_p\) on the class set with a polynomial in \(T_p\) on modular forms; the explicit factor is the constant \((p^{2k-2}-p)/(p-1)\).

In rank \(16\), \(\theta_{E_8 \perp E_8} = \theta_{D_{16}^+}\) by Problem~\ref{sw-pr:rank16} of Chapter~\ref{sw-chap}.

\begin{problem}\label{pr:kernel16}
The eigenvector of Problem~\ref{pr:rank16} with eigenvalue \(1800\) lies in the kernel of \(\theta\).  For the constant function, the eigenvalue \(32895\) is consistent with Problem~\ref{pr:commute} at \(p = 2\): \(129^2 - 2^7 + (2^{14} - 2) = 32895\).
\end{problem}

The two lattices are distinguished by their theta series of degree \(4\), and not before \cite{CL}.  A lift that kills an eigenvector says nothing about its eigenvalue; one needs a lift to a larger group, as in the Rallis inner product formula (Section~\ref{dn-sec:norm}).


\section{Hermitian lattices of rank five}\label{sec:hermitian}

Let \(\Oc = \Z[\ii]\), let \(p \equiv 1 \pmod 4\) be a prime, and let \(p = \pi\bar\pi\) in \(\Oc\).  Let \(L\) be a unimodular Hermitian lattice of rank \(n\), as in Chapter~\ref{dt-chap}, and let \(L^\sharp\) denote duals.  Then \(L/\pi L\) is an \(n\)-dimensional vector space over \(\Oc/\pi = \F_p\).  For a hyperplane \(H \subset L/\pi L\) let \(L_H\) be its preimage in \(L\), and
\[
  L'_H := L_H + \pi L_H^\sharp .
\]

\begin{problem}\label{pr:herm-nb}
\(L'_H\) is unimodular and lies in the genus of \(L\), with \(L \cap L'_H = L_H\) of index \(p\) in both.  Moreover \(L'_H \subset \bar\pi^{-1}L\), and \(L\) is obtained from \(L'_H\) by the same construction with \(\bar\pi\) in place of \(\pi\).
\end{problem}

We call \(L'_H\) a \emph{\(\pi\)-neighbour} of \(L\).  There are \((p^n - 1)/(p-1)\) of them, and \(T_\pi\varphi(L) := \sum_H \varphi(L'_H)\) defines an operator on functions on \(\Cl_n\).  By Problem~\ref{pr:herm-nb} and the argument of Problem~\ref{pr:nb-adjoint}, \(T_{\bar\pi}\) is the adjoint of \(T_\pi\).  Let \(r_L(m) := \#\{x \in L : h(x,x) = m\}\) and \(\theta_L := \sum_m r_L(m)\, q^m\), a modular form of weight \(n\) for \(\Gamma_0(4)\) with character \(\chi^n\), where \(\chi\) is the character modulo \(4\) \cite[Ch.~10]{Iw}.  Since \(\chi(p) = 1\), \(T_p\) acts on such forms by \(a(m) \mapsto a(pm) + p^{n-1}a(m/p)\).

\begin{problem}\label{pr:herm-count}
For \(x \in L\), the number of \(H\) with \(x \in L'_H\) is \((p^{n-1}-1)/(p-1)\) if \(x \notin \pi L\), and \((p^n - 1)/(p - 1)\) if \(x \in \pi L\).  A vector \(x \notin L\) lies in some \(L'_H\) exactly when \(y := \bar\pi x \in L \setminus \bar\pi L\) and \(p \mid h(y,y)\), and then in exactly one.
\end{problem}

\begin{problem}\label{pr:herm-eichler}
\(\displaystyle \sum_H \theta_{L'_H} = T_p\theta_L + \frac{p^{n-1} - p}{p - 1}\,\theta_L\).  Hence, with \(\theta(\varphi) := \sum_{L \in \Cl_n} \varphi(L)\theta_L/\lvert\Gamma_L\rvert\), one has \(\theta(T_\pi\varphi) = \bigl(T_p + \frac{p^{n-1}-p}{p-1}\bigr)\theta(\varphi)\), and likewise for \(\bar\pi\).
\end{problem}

Here the factor is linear in \(T_p\), because at a split prime the local group is \(\mathrm{GL}_n(\Q_p)\) and a \(\pi\)-neighbour is an edge of type one in the sense of Section~\ref{sec:split}.

Now take \(n = 5\), \(p = 5\) and \(\pi = 2 + \ii\).  By Problem~\ref{dt-pr:rank-five}, \(\Cl_5 = \{I_5, E \perp I_1\}\), with \(\lvert\Gamma\rvert = 122880\) and \(184320\), and \(r(1) = 20\) and \(4\).

\begin{problem}\label{pr:herm-brandt}
The matrix of \(T_\pi\) on \(\Cl_5\) is \(\begin{pmatrix} 525 & 256 \\ 384 & 397 \end{pmatrix}\), with eigenvalues \(781\) and \(141\).
\end{problem}

We checked this matrix by constructing all \(781\) neighbours of each lattice by machine.  Let \(g\) be the cusp form called \(f\) in Chapter~\ref{sw-chap}, 
\[
  g = \frac14\sum_{\alpha \in \Oc}\alpha^4 q^{\lvert\alpha\rvert^2} = \eta(\tau)^4\eta(2\tau)^2\eta(4\tau)^4 = q - 4q^2 + 16q^4 - 14q^5 - 64q^8 + 81q^9 + \cdots ;
\]
it spans the cusp forms of weight \(5\) for \(\Gamma_0(4)\) with character \(\chi\), and it has complex multiplication by \(\Q(\ii)\).  By Problem~\ref{sw-pr:cusp5}, \(\theta_{I_5} - \theta_{E \perp I_1} = 16\,g\).

\begin{problem}\label{pr:herm-theta}
The eigenvector of \(T_\pi\) with eigenvalue \(141\) is \(\varphi = (2, -3)\), and \(\theta(\varphi) = g/3840\).
\end{problem}

\begin{problem}\label{pr:herm-check}
\(a_5(g) = -14 = 2\operatorname{Re}\bigl((2+\ii)^4\bigr)\), and \(141 = -14 + 155\), as Problem~\ref{pr:herm-eichler} predicts.
\end{problem}

Here the lift is injective, and the Brandt matrix and the Hecke eigenvalues of \(\theta_{I_5}\) and \(\theta_{E \perp I_1}\) determine each other.


\section{Brandt matrices}\label{sec:brandt}

The original setting of Eichler's relation is a quaternion algebra, where the theta series of the norm forms of maximal orders give the modular forms of weight two and prime level \cite{Ei73}, \cite{Pi}, \cite{Gr}.  We follow Gross \cite[\S\S1--2]{Gr} in the smallest case with two classes.

Let \(B = (-1, -11)/\Q\), with \(\ii^2 = -1\), \(j^2 = -11\) and \(k = \ii j\); it is ramified exactly at \(11\) and \(\infty\).  Let \(\Oc_1 := \Z + \Z\ii + \Z\frac{1+j}{2} + \Z\frac{\ii + k}{2}\).  There is one other class of maximal orders, represented by an order \(\Oc_2\) with \(\Theta_{\Oc_2} := \sum_{x \in \Oc_2} q^{N(x)} = 1 + 6q + 6q^3 + 24q^4 + 18q^5 + \cdots\).  Let \(w_i := \lvert\Oc_i^\times\rvert/2\).  Gross defines the Brandt matrices \(B(m)\), \(m \ge 1\), whose entries \(B_{ij}(m)\) count the elements of norm \(m\) in \(\mathrm{Hom}(I_j, I_i)\), divided by \(2w_j\), for left ideals \(I_i\) of a fixed maximal order with right orders \(\Oc_i\); he shows that \(B_{ii}(m) = r_{\Oc_i}(m)/2w_i\), that \(w_j B_{ij}(m) = w_i B_{ji}(m)\), that \(\sum_j B_{ij}(m) = \sigma(m)\) for \(11 \nmid m\), and that \(m \mapsto B(m)\) is the action of the Hecke operators \(T_m\) on \(M_2(\Gamma_0(11))\) \cite[\S\S1--2]{Gr}.  For a prime \(\ell \ne 11\), \(B_{ij}(\ell)\) counts the left ideals \(J \subset I_i\) of the fixed order with \([I_i : J] = \ell^2\) and \(J \cong I_j\); there are \(\ell + 1\) such \(J\) in all, and they play the role of the \(\ell\)-neighbours of Section~\ref{sec:neighbours}.

\begin{problem}\label{pr:order}
\(\Oc_1\) is an order, its norm form is \((a^2 + ac + 3c^2) + (b^2 + bd + 3d^2)\), and \(\Oc_1^\times = \{\pm1, \pm\ii\}\).  Thus \(\Theta_{\Oc_1} = \bigl(\sum_{a,c} q^{a^2+ac+3c^2}\bigr)^2\).
\end{problem}

Eichler's mass formula \(\sum_i 1/w_i = (11 - 1)/12\) then gives \(w_2 = 3\) \cite{Gr}.

\begin{problem}\label{pr:brandt}
\(B(2) = \begin{pmatrix} 1 & 2 \\ 3 & 0 \end{pmatrix}\) and \(B(3) = \begin{pmatrix} 2 & 2 \\ 3 & 1 \end{pmatrix}\), with eigenvalues \(3, -2\) and \(4, -1\).
\end{problem}

\begin{problem}\label{pr:basis}
Let \(f := q\prod_{n \ge 1}(1 - q^n)^2(1 - q^{11n})^2 = q - 2q^2 - q^3 + 2q^4 + q^5 + \cdots\).  Then \(\Theta_{\Oc_1} - \Theta_{\Oc_2} = -2f\), and \((1, -1)\,B(\ell) = a_\ell(f)\,(1, -1)\) for \(\ell = 2, 3\).
\end{problem}

The eigenvalues \(-2\) and \(-1\) are the traces of Frobenius of the elliptic curve \(X_0(11)\) at \(2\) and \(3\).


\section{A split place}\label{sec:split}

In the chain of lifts of Adriazola's reconstruction \cite[\S10]{Adr}, the kernel of a lift at a split place is the indicator of \(xL \subset L'\) for a matrix \(x\), and the Hecke operators are sums over edges of lattices.  The relation between them is a count of subspaces, of which Problems~\ref{pr:herm-count} and~\ref{pr:herm-eichler} are the case of rank one.

Let \(Q\) be a power of \(p\), and let \(W\) be the unramified extension of \(\Z_p\) with residue field \(\F_Q\).  For lattices \(L'' \subset L\) in a \(W[1/p]\)-vector space, \((L'', L)\) is an \emph{edge of type \(i\)} if \(pL \subset L'' \subset L\) and \(\dim L/L'' = i\).  Let \(L \subset W[1/p]^N\) and \(L' \subset W[1/p]^{N+g}\) be lattices, let \(x\) be an \((N+g) \times N\) matrix with \(xL \subset L'\), and let \(\rho\) be the rank of \(x\) modulo \(p\) as a map \(L/pL \to L'/pL'\).

\begin{problem}\label{pr:split-count}
The number of edges \((L_1, L')\) of type \(i\) with \(xL \subset L_1\) is \(\qbinom{N+g-\rho}{i}{Q}\).  The number of edges \((L, A_1)\) of type \(i\) with \(xA_1 \subset L'\) is \(\qbinom{N-\rho}{i}{Q}\).
\end{problem}

\begin{problem}\label{pr:split-vandermonde}
As functions of \(x\) (with \(\mathbf 1_{L,L'}(x) := [xL \subset L']\)),
\[
  \sum_{(L_1, L') \text{ of type } i} \mathbf 1_{L, L_1}
  = \sum_{l=0}^{i} Q^{(i-l)(g-l)} \qbinom{g}{l}{Q} \sum_{(L, A_1) \text{ of type } i-l} \mathbf 1_{A_1, L'} .
\]
For \(N = 1\), \(i = 1\) and \(Q = p\) this is the first count of Problem~\ref{pr:herm-count}.
\end{problem}

\begin{problem}\label{pr:split-el}
For \(0 \le l \le g\), \(e_l(Q^{t/2}, Q^{t/2 - 1}, \dots, Q^{-t/2}) = Q^{-l(g-l)/2}\qbinom{g}{l}{Q}\), where \(t = g - 1\), the arguments step by \(Q^{-1}\), and \(e_l\) is the elementary symmetric polynomial.
\end{problem}

Summing Problem~\ref{pr:split-vandermonde} against a kernel that depends on \(x\) only through \(\mathbf 1_{L,L'}(x)\) and on the lengths of the lattices through a power of \(p\), and moving the sum over edges at the source to the other variable as in Problem~\ref{pr:nb-adjoint}, gives a relation of the form \(E_{j+1}(X)\, K = K\, P(X)\, E_j(uX)\) between the generating series \(E(X) = \sum_i \mathcal E_i X^i\) of normalized Hecke operators at the target and at the source.  Here \(u\) is a unit and \(P(X) = \prod_a (1 + aX)\) is scalar, with roots read off from Problem~\ref{pr:split-el}.  This is the relation for one lift in \cite[\S10.4]{Adr}.

\begin{problem}\label{pr:split-chain}
Suppose that \(E_{j+1}(X)\,K_j = K_j\, P_{j+1}(X)\, E_j(X)\) modulo \(X^{r+1}\) for \(0 \le j < m\), with \(P_{j+1}(X) = \prod_{a \in x_{j+1}}(1 + aX)\), and where \(E_j(X) = \sum_i \mathcal E^{(j)}_i X^i\), and let \(Y\) be the union of the families \(x_1, \dots, x_m\) and \(\mathcal K := K_{m-1}\cdots K_0\).  Then
\(\mathcal K\, \mathcal E^{(0)}_r = \sum_{l=0}^r (-1)^l h_l(Y)\, \mathcal E^{(m)}_{r-l}\, \mathcal K\),
where \(h_l\) is the complete symmetric polynomial.
\end{problem}

So a Hecke operator at the source of a chain of lifts moves to the target as a polynomial in the Hecke operators there, with coefficients \(h_l(Y)\).  Problem~\ref{pr:commute} is an instance with one lift.  For divisibility along a chain, the valuations of these coefficients are what matter \cite[\S10.2]{Adr}.

\section*{Hints}

\subsection*{\S\ref{sec:neighbours}}
\begin{description}
\item[\ref{pr:nb-line}] \(v/p\) has integral products with \(L_v\), and \((v/p) \cdot (v/p) \in 2\Z\).  Compare discriminants: \([L : L_v] = p\).  Conversely \(M \cap L = L_v\) for any \(v \in pM \setminus pL\).  Changing \(v\) by \(pL_v\) or by a unit does not change \(M_v\); to lift an isotropic \(\bar v\) to \(v \cdot v \equiv 0 \pmod{2p^2}\), adjust by \(p w\) with \(w \cdot v \not\equiv 0\).
\item[\ref{pr:nb-count}] The discriminant of \(L/pL\) is \((-1)^k\) times a square.  A split space of dimension \(2k\) has \(p^{2k-1} + p^k - p^{k-1}\) isotropic vectors, counting \(0\).
\item[\ref{pr:nb-adjoint}] Count the \(\Gamma\)-orbits of pairs \((L, M)\) with \(M \in \Nb_p(L)\), weighted by the inverse order of the stabilizer of the pair, in two ways.  The relation is symmetric because \(L \in \Nb_p(M)\).
\item[\ref{pr:rank16}] Row sums are \(32895\); use Problem~\ref{pr:nb-adjoint} for the second row.
\end{description}

\subsection*{\S\ref{sec:eichler}}
\begin{description}
\item[\ref{pr:count-in}] By Problem~\ref{pr:nb-line}, \(x \in M_v\) exactly when \(\bar x \perp \bar v\).  Count isotropic lines in \(\bar x^\perp\): if \(\bar x\) is anisotropic, \(\bar x^\perp\) is non-degenerate of odd dimension \(2k-1\); if \(\bar x\) is isotropic, \(\bar x^\perp = \F_p\bar x \oplus W\) with \(W\) split of dimension \(2k-2\).
\item[\ref{pr:count-out}] Use the description of \(M_v\) in Problem~\ref{pr:nb-line}.
\item[\ref{pr:eichler-check}] \(\sigma_3(9) = 757\) and \(\sigma_3(4) = 73\).
\item[\ref{pr:commute}] Exchange the sums by Problem~\ref{pr:nb-adjoint}, then use Theorem~\ref{thm:eichler}.  On an eigenform, \(T_{p^2} = T_p^2 - p^{k-1}\).
\item[\ref{pr:kernel16}] \(M_8(\mathrm{SL}_2(\Z))\) is spanned by \(E_4^2\), and \(T_2 E_4^2 = (1 + 2^7)E_4^2\).  The eigenvector is orthogonal to the constants.
\end{description}

\subsection*{\S\ref{sec:hermitian}}
\begin{description}
\item[\ref{pr:herm-nb}] Work locally.  Over \(\Z_p\), \(\Oc \otimes \Z_p = \Z_p \times \Z_p\), and a self-dual lattice is \(\Lambda \oplus \Lambda^\vee\) for a \(\Z_p\)-lattice \(\Lambda\).
\item[\ref{pr:herm-count}] Use the local description of \(L'_H\) from Problem~\ref{pr:herm-nb}, one component at a time.
\item[\ref{pr:herm-eichler}] Sum Problem~\ref{pr:herm-count} over the vectors of norm \(m\), using \(h(\bar\pi y, \bar\pi y) = p\, h(y,y)\) and \(r_{\pi L}(m) = r_L(m/p)\).
\item[\ref{pr:herm-brandt}] Compare the coefficients of \(q\) in Problem~\ref{pr:herm-eichler} for \(L = I_5\), using \(r_{I_5}(5) = 8424\) and the values of \(r(1)\).  For the second row, do the same for \(L = E \perp I_1\), with \(r_{E \perp I_1}(5) = 8648\); or use Problem~\ref{pr:nb-adjoint} together with complex conjugation, which exchanges \(\pi\)- and \(\bar\pi\)-neighbours and fixes both classes.
\item[\ref{pr:herm-theta}] Use \(122880 = 2 \cdot 61440\) and \(184320 = 3 \cdot 61440\).
\item[\ref{pr:herm-check}] Compute \(a_5(g)\) from the product.  Then apply Problem~\ref{pr:herm-eichler} to \(\varphi\) and compare with Problem~\ref{pr:herm-theta}.
\end{description}

\subsection*{\S\ref{sec:brandt}}
\begin{description}
\item[\ref{pr:order}] Write an element as \((a + c/2) + (b + d/2)\ii + (c/2)j + (d/2)k\).
\item[\ref{pr:brandt}] Find the diagonal entries from \(r_{\Oc_1}(m)\) and from \(\Theta_{\Oc_2}\), and the others from the row sums.  Check \(w_2B_{12} = w_1B_{21}\).
\item[\ref{pr:basis}] Both sides lie in \(M_2(\Gamma_0(11))\), which has dimension \(2\); compare the coefficients of \(q^0, q^1, q^2\).
\end{description}

\subsection*{\S\ref{sec:split}}
\begin{description}
\item[\ref{pr:split-count}] \(L_1 \supset pL'\) corresponds to a subspace of \(L'/pL'\) of codimension \(i\) containing the image of \(x\).  \(A_1 \subset p^{-1}L\) corresponds to the subspace \(pA_1/pL\) of \(L/pL\), which must lie in the kernel of \(x\) modulo \(p\).
\item[\ref{pr:split-vandermonde}] Problem~\ref{pr:split-count} and the \(q\)-Vandermonde identity \cite[Ch.~3]{Andr}\[\qbinom{a+b}{i}{Q} = \sum_l Q^{(i-l)(b-l)}\qbinom{b}{l}{Q}\qbinom{a}{i-l}{Q} .\]
\item[\ref{pr:split-el}] Use the Gaussian binomial theorem \[\prod_{s=0}^{g-1}(1 + Q^s z) = \sum_l Q^{l(l-1)/2}\qbinom{g}{l}{Q} z^l .\]
\item[\ref{pr:split-chain}] Work with power series in \(X\) modulo \(X^{r+1}\); the scalars commute with the operators, and \(\prod_{a \in Y}(1 + aX)^{-1} = \sum_l (-1)^l h_l(Y) X^l\).
\end{description}

{\small\begin{thebibliography}{Andr}
\bibitem[Adr]{Adr} J.~Adriazola, \emph{The Alp\"oge proof of the Riemann hypothesis}, manuscript, revised 24 November 2026, \url{https://github.com/jimjamzola/RH-Alpoge-Reconstruction}.
\bibitem[An]{An} A.~N.~Andrianov, \emph{Quadratic Forms and Hecke Operators}, Grundlehren Math.\ Wiss.\ 286, Springer, 1987.
\bibitem[Andr]{Andr} G.~E.~Andrews, \emph{The Theory of Partitions}, Encyclopedia Math.\ Appl.\ 2, Addison-Wesley, 1976.
\bibitem[CL]{CL} G.~Chenevier and J.~Lannes, \emph{Automorphic Forms and Even Unimodular Lattices}, Ergeb.\ Math.\ Grenzgeb.\ (3) 69, Springer, 2019.
\bibitem[CoS]{CoS} J.~H.~Conway and N.~J.~A.~Sloane, \emph{Sphere Packings, Lattices and Groups}, 3rd ed., Grundlehren Math.\ Wiss.\ 290, Springer, 1999.
\bibitem[Ei52]{Ei52} M.~Eichler, \emph{Quadratische Formen und orthogonale Gruppen}, Grundlehren Math.\ Wiss.\ 63, Springer, 1952.
\bibitem[Ei73]{Ei73} M.~Eichler, \emph{The basis problem for modular forms and the traces of the Hecke operators}, in \emph{Modular Functions of One Variable I}, Lecture Notes in Math.\ 320, Springer, 1973, 75--151.
\bibitem[Gr]{Gr} B.~H.~Gross, \emph{Heights and the special values of \(L\)-series}, in \emph{Number Theory (Montreal, 1985)}, CMS Conf.\ Proc.\ 7, Amer.\ Math.\ Soc., 1987, 115--187.
\bibitem[Iw]{Iw} H.~Iwaniec, \emph{Topics in Classical Automorphic Forms}, Grad.\ Stud.\ Math.\ 17, Amer.\ Math.\ Soc., 1997.
\bibitem[Kn]{Kn} M.~Kneser, \emph{Klassenzahlen definiter quadratischer Formen}, Arch.\ Math.\ \textbf{8} (1957), 241--250.
\bibitem[Pi]{Pi} A.~Pizer, \emph{An algorithm for computing modular forms on \(\Gamma_0(N)\)}, J.\ Algebra \textbf{64} (1980), 340--390.
\bibitem[Ra]{Ra} S.~Rallis, \emph{Langlands' functoriality and the Weil representation}, Amer.\ J.\ Math.\ \textbf{104} (1982), 469--515.
\bibitem[SeA]{SeA} J.-P.~Serre, \emph{A Course in Arithmetic}, Grad.\ Texts in Math.\ 7, Springer, 1973.
\end{thebibliography}}

\end{document}
